\section{Polygon skein algebras and Hall-compatible quantization}
\label{sec:skein}

We now fix the target algebra.  The purpose of this section is not to reprove
the skein--cluster theorem, but to state precisely the presentation and basis
that will be compared with the Hall algebra.

\subsection{Marked polygon and coefficient system}

Let $P=P_{n+3}$ be an oriented convex polygon with marked vertices
$0,1,\ldots,n+2$ in counterclockwise order.  A boundary edge is regarded as a
frozen arc, and an internal diagonal as a mutable arc.  Fix a triangulation
$T$ and the corresponding initial cluster seed.

The previous Hall-globalization theorem supplies a quantum commutation form on
each rigid chamber.  We choose a geometric coefficient system whose quantum
seed has exactly these forms.  There are two standard regimes.

\begin{enumerate}[label=(\roman*),leftmargin=2.3em]
\item If the mutable exchange matrix is full rank, one may use the
      coefficient-free compatible form and then adjoin/invert boundary arcs.
\item In arbitrary rank, use the principal Hall-compatible parameter and its
      walled-surface realization.  The walls encode the frozen coefficient
      monomials, while the mutable chamber forms are those pulled back from the
      Hall term lattice.
\end{enumerate}

The existence of skein models with arbitrary geometric coefficients for
unpunctured surfaces is proved by Ishibashi--Kano--Yuasa
\cite{IshibashiKanoYuasa}.  Since the theorem below only uses the local
crossing relation, the non-crossing basis and the equality with the quantum
cluster algebra, either convention may be used.  We denote the resulting
frozen-localized skein algebra by
\[
 \Skein_\omega^{\rm fr}(P)
\]
and its quantum cluster algebra by
\[
 \Aq=\mathcal A_\omega^{\rm fr}(P).
\]
The coefficient ring is $\mathbb Q(\omega)$, and all boundary arcs are
invertible.

\subsection{Local skein relations}

For an arc or a simple multicurve $D$, write $s_D$ for its skein class.  The
product is vertical superposition.  At a crossing of two arcs there are two
smoothings, denoted
\[
 \gamma\#_0\delta,
 \qquad
 \gamma\#_\infty\delta.
\]
Each is a two-component multicurve; a component may be a boundary edge.
Endpoint elevations and wall states are ordered by the fixed convention of the
coefficient system.

After the usual normalization of stated arcs, the local relation takes the
form
\begin{equation}
 s_\gamma s_\delta
 =\omega^{a(\gamma,\delta)}s_{\gamma\#_0\delta}
  +\omega^{b(\gamma,\delta)}s_{\gamma\#_\infty\delta}.
 \label{eq:quantum-skein-relation}
\end{equation}
For the standard Kauffman-bracket normalization one may take
$a=1$, $b=-1$ for one ordering and interchange the exponents for the reverse
ordering.  In a walled coefficient model, frozen Laurent monomials can be
absorbed into the two smoothing symbols.  We retain the symbols $a,b$ because
the Hall quotient theorem is independent of this normalization.

Two non-crossing arcs quasi-commute:
\begin{equation}
 s_\gamma s_\delta
 =\omega^{\lambda_T(\gamma,\delta)}s_\delta s_\gamma.
 \label{eq:skein-quasicommutation}
\end{equation}
If they belong to one triangulation, the exponent is the corresponding entry
of the quantum commutation matrix of that seed.  This is precisely the form
already reproduced by the Hall cocycle in
\eqref{eq:compatible-Hall-quantum}.

The contractible-loop and boundary-state relations are part of the standard
skein presentation.  In a polygon, every closed contractible component is
removed by a scalar, and every boundary edge is a frozen unit after
localization.  These relations match the central specialization in
\cref{def:Hred}.

\subsection{The non-crossing basis}

A simple stated multicurve in $P$ is a finite union of pairwise non-crossing
arcs, with the prescribed endpoint order and with no removable contractible
component.  Multiplicities are allowed.  Let
\[
 \mathfrak N(P)
\]
denote their isotopy classes after boundary edges have been recorded as frozen
Laurent monomials.

\begin{theorem}[Polygon skein basis]
\label{thm:skein-basis}
The elements
\[
 \{s_D:D\in\mathfrak N(P)\}
\]
form a $\mathbb Q(\omega)$-basis of
$\Skein_\omega^{\rm fr}(P)$.  Every crossing diagram has a finite and
choice-independent expansion in this basis obtained by repeated use of
\eqref{eq:quantum-skein-relation}.
\end{theorem}

\begin{proof}[Source of the theorem]
For marked surfaces, Muller proves the multicurve basis and the skein--cluster
inclusions, with equality after Ore localization when each component has at
least two marked points \cite{Muller}.  The polygon satisfies this hypothesis.
The walled coefficient extension is established in
\cite{IshibashiKanoYuasa}.  For stated skein algebras, the polygon equality and
basis refinements are proved in \cite{CaoHuangWang}.  These results include the
local confluence needed here.
\end{proof}

The basis may also be understood by a direct Diamond-lemma argument: resolve a
chosen crossing, note that both smoothings reduce the total number of
crossings, and check the finitely many local overlap types.  The deterministic
code supplied with this paper implements this recursion for polygon chord
multidiagrams.

\subsection{Equality with the quantum cluster algebra}

The diagonal $\gamma$ is a quantum cluster variable, the boundary edges are
frozen variables, and a non-crossing multidiagonal is a normalized quantum
cluster monomial.  The surface skein--cluster theorem gives the following.

\begin{theorem}[Skein--cluster equality for polygons]
\label{thm:skein-cluster}
There is an algebra isomorphism
\begin{equation}
 \Psi_{\rm sc}:
 \Skein_\omega^{\rm fr}(P)
 \xrightarrow{\sim}
 \Aq
 \label{eq:skein-cluster-isomorphism}
\end{equation}
which sends every diagonal skein class to the correspondingly labelled
quantum cluster variable and every boundary edge to its frozen variable.
Moreover, for $D\in\mathfrak N(P)$,
\begin{equation}
 \Psi_{\rm sc}(s_D)
 =M_R(\mult_D)
 =\Theta_{g(D)}
 \label{eq:noncrossing-theta}
\end{equation}
for any cluster $R$ containing the support of $D$.  Here $M_R(\mult_D)$ is the
normalized based monomial and $\Theta_{g(D)}$ is the quantum theta function at
its cluster-chamber label.
\end{theorem}

\begin{proof}[Source of the theorem]
Muller proves equality of the localized skein, quantum cluster and quantum
upper cluster algebras for marked surfaces satisfying the boundary
hypothesis \cite{Muller}.  The coefficient and stated refinements used here
are supplied by \cite{IshibashiKanoYuasa,CaoHuangWang}.  The identification of
theta functions with normalized cluster monomials inside a cluster chamber is
also a consequence of the quantum theta theory of Davison--Mandel
\cite{DavisonMandel}.
\end{proof}

\subsection{Presentation by arcs}

Let $\mathcal F_P$ be the free associative algebra over
$\mathbb Q(\omega)$ generated by symbols $x_\gamma$ for all diagonals and
invertible symbols for the boundary edges.  Impose:
\begin{enumerate}[label=(\roman*),leftmargin=2.3em]
\item the boundary and contractible-loop relations;
\item the quasi-commutation relations for non-crossing arcs;
\item the two-smoothing relation \eqref{eq:quantum-skein-relation} for every
      ordered crossing pair.
\end{enumerate}
Then \cref{thm:skein-basis} gives
\begin{equation}
 \mathcal F_P/\mathcal I_{\rm sk}
 \xrightarrow{\sim}
 \Skein_\omega^{\rm fr}(P),
 \label{eq:skein-presentation}
\end{equation}
where $\mathcal I_{\rm sk}$ is the ideal generated by these local relations.
The irreducible normal words are indexed by $\mathfrak N(P)$.

The Hall algebra already satisfies the compatible quasi-commutations and the
boundary relations after the cocycle twist and central reduction.  Thus only
the crossing relations have to be imposed on $\Hred$.
